
% =========================
%  Appendices - Vehicle ReID
% =========================
\appendix

\section{Исходный код пайплайна}
\label{app:code}
Ключевые фрагменты кода из Jupyter notebooks обучения.

\clearpage
\subsection{Модель ResNet-50 + BNNeck}
\begin{lstlisting}[language=Python, caption={BNNeck и ResNet-50}]
class BNNeck(nn.Module):
    def __init__(self, in_features, num_classes):
        super().__init__()
        self.bn = nn.BatchNorm1d(in_features)
        self.bn.bias.requires_grad_(False)
        self.classifier = nn.Linear(in_features, num_classes, bias=False)
        nn.init.normal_(self.bn.weight, 1.0, 0.02)
        nn.init.normal_(self.classifier.weight, std=0.001)

    def forward(self, x):
        bn_x = self.bn(x)
        logits = self.classifier(bn_x)
        return bn_x, logits

class ResNet50ReID(nn.Module):
    def __init__(self, num_classes, pretrained=True):
        super().__init__()
        resnet = models.resnet50(
            weights='IMAGENET1K_V1' if pretrained else None)
        resnet.layer4[0].downsample[0].stride = (1, 1)
        resnet.layer4[0].conv2.stride = (1, 1)
        self.backbone = nn.Sequential(
            resnet.conv1, resnet.bn1, resnet.relu, resnet.maxpool,
            resnet.layer1, resnet.layer2, resnet.layer3, resnet.layer4)
        self.gap = nn.AdaptiveAvgPool2d(1)
        self.bnneck = BNNeck(2048, num_classes)

    def forward(self, x):
        feat_map = self.backbone(x)
        features = self.gap(feat_map).flatten(1)
        bn_features, logits = self.bnneck(features)
        if self.training:
            return features, bn_features, logits
        return F.normalize(bn_features, p=2, dim=1)
\end{lstlisting}

\clearpage
\subsection{Модель ViT-Small + BNNeck}
\begin{lstlisting}[language=Python, caption={ViT-Small с timm}]
import timm

class ViTReID(nn.Module):
    def __init__(self, num_classes, model_size='small', pretrained=True):
        super().__init__()
        model_names = {
            'tiny': 'vit_tiny_patch16_224',
            'small': 'vit_small_patch16_224',
            'base': 'vit_base_patch16_224'}
        self.backbone = timm.create_model(
            model_names[model_size], pretrained=pretrained, num_classes=0)
        self.embed_dim = self.backbone.embed_dim
        self.bnneck = BNNeck(self.embed_dim, num_classes)

    def forward(self, x):
        features = self.backbone(x)
        bn_features, logits = self.bnneck(features)
        if self.training:
            return features, bn_features, logits
        return F.normalize(bn_features, p=2, dim=1)
\end{lstlisting}

\clearpage
\subsection{Функции потерь}
\begin{lstlisting}[language=Python, caption={Cross-Entropy и Triplet Loss}]
class CrossEntropyLabelSmooth(nn.Module):
    def __init__(self, num_classes, epsilon=0.1):
        super().__init__()
        self.num_classes = num_classes
        self.epsilon = epsilon
        self.logsoftmax = nn.LogSoftmax(dim=1)

    def forward(self, inputs, targets):
        log_probs = self.logsoftmax(inputs)
        targets_one_hot = torch.zeros_like(log_probs).scatter_(
            1, targets.unsqueeze(1), 1)
        targets_smooth = ((1 - self.epsilon) * targets_one_hot
            + self.epsilon / self.num_classes)
        loss = (-targets_smooth * log_probs).sum(dim=1).mean()
        return loss

class TripletLoss(nn.Module):
    def __init__(self, margin=0.3):
        super().__init__()
        self.margin = margin

    def forward(self, features, labels):
        dist_mat = torch.cdist(features, features, p=2)
        mask_pos = labels.unsqueeze(0) == labels.unsqueeze(1)
        mask_neg = ~mask_pos
        dist_pos = dist_mat.clone()
        dist_pos[~mask_pos] = 0
        hard_pos = dist_pos.max(dim=1)[0]
        dist_neg = dist_mat.clone()
        dist_neg[~mask_neg] = float('inf')
        hard_neg = dist_neg.min(dim=1)[0]
        loss = F.relu(hard_pos - hard_neg + self.margin).mean()
        return loss
\end{lstlisting}

\clearpage
\subsection{PK-семплер}
\begin{lstlisting}[language=Python, caption={PK-Sampler для батчевого семплинга}]
class PKSampler:
    def __init__(self, dataset, p, k):
        self.p, self.k = p, k
        self.idx_to_samples = dataset.idx_to_samples
        self.labels = list(self.idx_to_samples.keys())

    def __iter__(self):
        random.shuffle(self.labels)
        batch = []
        for label in self.labels:
            indices = self.idx_to_samples[label]
            if len(indices) >= self.k:
                selected = random.sample(indices, self.k)
            else:
                selected = random.choices(indices, k=self.k)
            batch.extend(selected)
            if len(batch) >= self.p * self.k:
                yield batch[:self.p * self.k]
                batch = batch[self.p * self.k:]

    def __len__(self):
        return len(self.labels) // self.p
\end{lstlisting}

\clearpage
\subsection{Датасет VeRi-776}
\begin{lstlisting}[language=Python, caption={Загрузчик датасета}]
class VeRiDataset(Dataset):
    def __init__(self, data_dir, split='train', transform=None):
        self.data_dir = data_dir
        self.transform = transform
        if split == 'train':
            self.img_dir = os.path.join(data_dir, 'image_train')
        elif split == 'query':
            self.img_dir = os.path.join(data_dir, 'image_query')
        else:
            self.img_dir = os.path.join(data_dir, 'image_test')

        self.samples = []
        self.label_to_idx = {}
        self.idx_to_samples = defaultdict(list)

        filenames = [f for f in os.listdir(self.img_dir)
                     if f.endswith('.jpg')]
        for fname in filenames:
            parts = fname.split('_')
            vid = int(parts[0])
            cam = int(parts[1][1:]) if parts[1].startswith('c') else 0
            if vid not in self.label_to_idx:
                self.label_to_idx[vid] = len(self.label_to_idx)
            label = self.label_to_idx[vid]
            self.samples.append((fname, label, cam))
            self.idx_to_samples[label].append(len(self.samples) - 1)
        self.num_classes = len(self.label_to_idx)

    def __getitem__(self, idx):
        fname, label, cam = self.samples[idx]
        img = Image.open(os.path.join(self.img_dir, fname)).convert('RGB')
        if self.transform:
            img = self.transform(img)
        return {'image': img, 'label': label, 'camera_id': cam}
\end{lstlisting}

\clearpage
\subsection{Функции оценки}
\begin{lstlisting}[language=Python, caption={Вычисление метрик CMC и mAP}]
def compute_metrics(query_feats, gallery_feats,
                    query_labels, gallery_labels,
                    query_cams, gallery_cams):
    dist_mat = 1 - np.dot(query_feats, gallery_feats.T)
    num_query = len(query_labels)
    all_AP, all_cmc = [], np.zeros(50)

    for i in range(num_query):
        q_label, q_cam = query_labels[i], query_cams[i]
        valid_mask = ~((gallery_labels == q_label) &
                       (gallery_cams == q_cam))
        if valid_mask.sum() == 0:
            continue
        distances = dist_mat[i][valid_mask]
        g_labels = gallery_labels[valid_mask]
        indices = np.argsort(distances)
        matches = (g_labels[indices] == q_label).astype(np.int32)
        cmc = matches.cumsum()
        cmc[cmc > 1] = 1
        all_cmc[:len(cmc)] += cmc[:50]
        num_rel = matches.sum()
        if num_rel > 0:
            precision = matches.cumsum() / (np.arange(len(matches)) + 1)
            all_AP.append((precision * matches).sum() / num_rel)

    cmc = all_cmc / num_query * 100
    return {
        'Rank-1': cmc[0], 'Rank-5': cmc[4],
        'Rank-10': cmc[9], 'mAP': np.mean(all_AP) * 100}
\end{lstlisting}

\textbf{Примечание о junk-изображениях.} Протокол VeRi-776 определяет как «junk» все gallery-изображения той же камеры, что и запрос (файл \texttt{jk\_image.txt}). В приведённом листинге \texttt{valid\_mask} исключает из ранжирования изображения с \emph{совпадающим ID и той же камерой} (прямые дубли запроса). Изображения других транспортных средств той же камеры остаются в галерее как дополнительные негативные примеры (дистракторы), что делает оценку строже, а не мягче: модели сложнее найти правильный ответ в присутствии камерно-схожих негативов. В производственном коде (\texttt{vehicle\_reid/evaluation.py}) используется аналогичная логика исключения same-camera same-ID изображений.

\section{Структура проекта}
\label{app:structure}

\begin{verbatim}
code/
+-- train_resnet50.ipynb   # Jupyter notebook ResNet-50
+-- train_vit.ipynb        # Jupyter notebook ViT-Small
+-- comparison.ipynb       # Сравнение моделей
+-- vehicle_reid/
|   +-- model.py           # ResNet-50 + BNNeck
|   +-- vit_model.py       # ViT-Small + BNNeck
|   +-- losses.py          # CE + Triplet Loss
|   +-- dataset.py         # VeRi-776 loader
|   +-- transforms.py      # Augmentations
|   +-- reranking.py       # K-reciprocal re-ranking
|   +-- evaluation.py      # CMC/mAP metrics
+-- configs/
|   +-- resnet50.yaml
|   +-- vit.yaml
+-- requirements.txt

images/
+-- comparison_metrics.png
+-- training_curves_comparison.png
\end{verbatim}

\section{Воспроизводимость}
\label{app:repro}
\begin{itemize}
  \item \textbf{Random seed}: 42 (NumPy, PyTorch)
  \item \textbf{Версии}: PyTorch 2.0, torchvision 0.15, timm 1.0
  \item \textbf{Hardware}: Google Colab T4 GPU (16GB VRAM)
  \item \textbf{Датасет}: VeRi-776 (Kaggle)
\end{itemize}

